% chapters/chapter2.tex
\section{ПРОЄКТУВАННЯ ПРОГРАМНОГО РІШЕННЯ}

% Розділ 2 згідно нових методичних вказівок (12-15 стор.):
% 2.1 Вибір методів, технологій та засобів реалізації
% 2.2 Математичне та алгоритмічне забезпечення
% 2.3 Загальна архітектура / структура системи
% 2.4 Структури даних
% 2.5 Механізми обміну даними та інтеграції системи
% 2.6 Інтерфейс користувача та взаємодія з системою
% 2.7 Висновки до розділу 2
%
% Конкретні версії бібліотек, середовища розробки та деталі реалізації
% виносяться до розділу 3.

% ─────────────────────────────────────────────────
\subsection{Вибір методів, технологій та засобів реалізації}

У підрозд.~1.5 сформульовано технічне завдання: розробити веб-симулятор
потокового обчислювача 16-точкового ШПФ з інтерактивною візуалізацією
руху токенів, числовою верифікацією та відкритим доступом. Нижче
обґрунтовується вибір архітектурних підходів і сімейств інструментів,
що покривають ці вимоги. Конкретні версії бібліотек та параметри
середовища розробки наводяться у підрозд.~3.1.

\subsubsection*{Вибір архітектурного підходу}

Розглянуто три альтернативних підходи до побудови симулятора.

\textbf{Варіант~А: настільний застосунок (Electron, Qt, JavaFX).}
Забезпечує повний доступ до ресурсів хост-системи та багатопоточність,
однак вимагає встановлення на кожну робочу станцію, різних збірок під
Windows, Linux, macOS та погіршує переносимість --- прямо суперечить
вимозі NF7 (робота без встановлення) і ускладнює дотримання NF4
(портативність).

\textbf{Варіант~Б: клієнт-серверний веб-застосунок.} Тонкий клієнт у
браузері виконує лише візуалізацію, а ядро симуляції розміщується на
сервері. Такий підхід вносить мережеву затримку, що несумісно з
вимогою NF2 (час відгуку $\le 100$~мс), потребує підтримки серверної
інфраструктури та ускладнює відтворюваність результатів поза межами
хостингу.

\textbf{Варіант~В: повністю клієнтський SPA (Single-Page Application).}
--- \emph{обрано як основу архітектури}. Усі обчислення виконуються
на стороні браузера; серверна частина відсутня як така. Завантаження
застосунку відбувається одноразово як статичний пакет (HTML+JS+CSS),
після чого робота не потребує мережі. Цей варіант одночасно задовольняє
NF4, NF5 (детермінованість локального виконання), NF7 (без встановлення)
і забезпечує низький час відгуку за рахунок відсутності мережевих
викликів у петлі симуляції.

\subsubsection*{Вибір парадигми реалізації UI та стану}

Симулятор поєднує три принципово різні підсистеми: анімована
графова візуалізація DFG, числове ядро симуляції та панель керування
і метрик. Для їх ефективної координації обрано поєднання
\emph{реактивного рендерингу} та \emph{централізованого
спостерігаваного сховища стану}.

Реактивний рендеринг (React-парадигма) природно відображає
стан симулятора на візуальні компоненти: зміна значення в одному
полі сховища автоматично ініціює перемалювання лише тих компонентів,
що підписані на це поле. Це усуває ручне управління DOM і вкладається
у вимогу NF3 (стабільні 60~fps при швидкості симуляції $\le 5\times$).

Централізоване сховище стану (pattern single source of truth)
гарантує, що всі панелі --- GraphView, SpectrumView, Metrics ---
бачать один і той самий «знімок» стану у поточному кадрі. Це
виключає класичні помилки неузгодженості між візуальними
представленнями однієї сутності.

Альтернатива у вигляді прямого маніпулювання DOM (jQuery-style)
була відкинута через лінійне зростання складності при додаванні
нових компонентів, а Flux-підходи з ред'юсерами (Redux) --- через
надмірну кількість шаблонного коду для базових сценаріїв.

\subsubsection*{Вибір методу візуалізації DFG}

Для реалізації інтерактивного графа потоку даних розглянуто три
підходи:

\begin{itemize}
  \item \textbf{Рендеринг на HTML-canvas.} Забезпечує максимальну
        продуктивність при великій кількості об'єктів, але вимагає
        ручної реалізації перетягування, виділення та спливаючих
        підказок, що значно збільшує обсяг коду для UI-інтеракцій.
  \item \textbf{SVG + ручне кодування.} Простіший у налагодженні,
        але лише для невеликих графів; розширюваність обмежена
        тим, що кожен новий тип вузла потребує окремого обробника
        подій.
  \item \textbf{Спеціалізована бібліотека графової візуалізації
        (React~Flow).} --- \emph{обрано}. Надає готові механізми
        перетягування вузлів, зума, анімованих ребер, плаваючих
        підказок, а також формальну модель \texttt{Node}/\texttt{Edge},
        що добре відображається на абстракцію DFG.
\end{itemize}

\subsubsection*{Шаблони проєктування}

У проєктуванні системи застосовано такі класичні шаблони:

\begin{itemize}
  \item \textit{Observer (Publisher/Subscriber)} --- для подій
        \texttt{onToken}, \texttt{onFire}, \texttt{onOutput}:
        ядро симуляції є видавцем, а UI-компоненти
        (Metrics, теплова карта, SpectrumView) --- підписниками
        без прямих посилань на ядро;
  \item \textit{State / Finite State Machine} --- для керування
        життєвим циклом симуляції: \texttt{IDLE} $\to$
        \texttt{RUNNING} $\to$ \texttt{PAUSED} $\to$ \texttt{DONE};
  \item \textit{Strategy} --- для реалізації різних типів вузлів
        (\texttt{source}, \texttt{dft4}, \texttt{twiddle},
        \texttt{sink}): поведінка вузла інкапсулюється у функції,
        обрану за полем \texttt{NodeKind};
  \item \textit{Factory} --- для генерації графа через функцію
        \texttt{generateDFT4x4()}, що повертає готовий об'єкт
        \texttt{Graph} на основі параметрів $N$ та \texttt{latency};
  \item \textit{Facade} --- клас \texttt{DataflowRuntime} приховує
        внутрішню реалізацію черг і планувальника за єдиним методом
        \texttt{tick(dt, maxFires)}.
\end{itemize}

\subsubsection*{Типізація та засоби забезпечення якості}

Для дотримання NF1 (точність $\varepsilon \le 10^{-9}$) і NF5
(детермінованість) проєктом передбачено статичну типізацію вихідного
коду. Це унеможливлює помилки передачі аргументів неправильного
типу (наприклад, числа замість комплексного об'єкта) на етапі
компіляції. Комплексні числа подаються як структура \{\,re, im\,\}
з операціями у подвійній точності IEEE~754, що забезпечує необхідний
запас точності відносно $\varepsilon$.

Підсумкова таблиця зв'язків між вимогами підрозд.~1.5 та обраними
архітектурними рішеннями наведена у табл.~\ref{tab:req_to_tech}.

\begin{table}[H]
  \caption{Зв'язок вимог і архітектурних рішень}
  \label{tab:req_to_tech}
  \centering
  \small
  \begin{tabular}{p{2.2cm}p{6.0cm}p{5.8cm}}
    \toprule
    Вимога & Архітектурне рішення & Обґрунтування \\
    \midrule
    NF1 (точність) & Статична типізація + IEEE 754 double & Виключення помилок типів, $\varepsilon_{\mathrm{mach}}{\approx}2.22{\cdot}10^{-16}$ \\
    NF2 (відгук $\le 100$~мс) & Клієнтський SPA без серверних викликів & Відсутність мережевої затримки \\
    NF3 (60~fps) & Реактивний рендеринг з підпискою на окремі поля & Перемальовування лише змінених компонентів \\
    NF4 (портативність) & Стандарт ES2020, Web APIs & Робота у Chrome/Firefox/Edge без polyfill \\
    NF5 (детермінованість) & Логічний час + статичний граф & Ідентичні входи $\Rightarrow$ ідентичний журнал \\
    NF6 (MIT/Apache) & Вибір лише OSS-бібліотек & React, React Flow, Zustand, TypeScript \\
    NF7 (без встановлення) & SPA у статичному хостингу & Одноразове завантаження бандлу \\
    NF8 (зручність) & React Flow + панелі Controls/Metrics & Двокліковий доступ до базових операцій \\
    \bottomrule
  \end{tabular}
\end{table}


% ─────────────────────────────────────────────────
\subsection{Математичне та алгоритмічне забезпечення}

У цьому підрозділі подано математичний апарат 16-точкового ШПФ
з основою~4, що становить інтелектуальне ядро симулятора, та формалізацію
обчислень у парадигмі потоку даних. Виведені формули та схеми алгоритму
безпосередньо визначають структуру графа (підрозд.~2.3) і структури
даних (підрозд.~2.4).

\subsubsection*{Дискретне перетворення Фур'є}

Дискретне перетворення Фур'є (ДПФ) скінченної комплекснозначної
послідовності $\{x_n\}_{n=0}^{N-1}$ визначається
як~\cite{nussbaumer1982,protsko2025}:

\begin{equation}
  X_k = \sum_{n=0}^{N-1} x_n \, e^{-j 2\pi n k / N}
       = \sum_{n=0}^{N-1} x_n \, W_N^{n k},
  \quad k = 0, 1, \ldots, N-1,
  \label{eq:dft}
\end{equation}

де $W_N = e^{-j 2\pi/N}$ --- поворотний множник (twiddle factor).
Пряме обчислення~(\ref{eq:dft}) потребує $N^2$ комплексних множень
і $N(N-1)$ комплексних додавань, тобто має складність $O(N^2)$, що
недостатньо для задач реального часу при великих $N$.

Клас швидких алгоритмів перетворення Фур'є (ШПФ) знижує складність
до $O(N \log_2 N)$ за рахунок факторизації довжини
$N = N_1 \cdot N_2$ та заміни одного $N$-точкового ДПФ на систему
$N_1$ коротших $N_2$-точкових та $N_2$ коротших $N_1$-точкових
перетворень~\cite{cooley1965,nussbaumer1982}. Для $N = 16$ природним
вибором є алгоритм з основою~4 (radix-4), у якому $N_1 = N_2 = 4$.

\subsubsection*{Двовимірна $4 \times 4$ декомпозиція для $N = 16$}

Для $N = 16$ часовий і частотний індекси подаються у вигляді:

\begin{equation}
  n = 4s + r, \quad k = 4k_2 + k_1, \quad r, s, k_1, k_2 \in \{0,1,2,3\}.
  \label{eq:idx_map}
\end{equation}

Вхідна послідовність $\{x(n)\}_{n=0}^{15}$ вкладається у матрицю
$\mathbf{A}$ розміром $4 \times 4$:

\begin{equation}
  A[r][s] = x(4s + r), \quad r, s \in \{0,1,2,3\}.
  \label{eq:mat_A}
\end{equation}

Підстановка~(\ref{eq:idx_map}) у~(\ref{eq:dft}) із використанням
тотожності $W_{16}^{4 \cdot (\cdot)} = W_4^{(\cdot)}$
дає~\cite{burrus1985dft,nussbaumer1982}:

\begin{equation}
  X(4k_2 + k_1)
    = \sum_{r=0}^{3} W_{16}^{r k_1}
      \underbrace{\Bigl( \sum_{s=0}^{3} A[r][s]\, W_4^{s k_1} \Bigr)}_{G_r(k_1)}
      W_4^{r k_2},
  \label{eq:fft4x4}
\end{equation}

що визначає тристадійну схему алгоритму.

\textbf{Стадія~1 --- рядкові 4-точкові ДПФ.}
Для кожного $r \in \{0,1,2,3\}$ обчислюється 4-точкове ДПФ рядка:

\begin{equation}
  G_r(k_1) = \sum_{s=0}^{3} A[r][s]\, W_4^{s k_1},
  \quad k_1 \in \{0,1,2,3\}.
  \label{eq:stage1}
\end{equation}

Чотири перетворення~(\ref{eq:stage1}) є взаємно незалежними та
виконуються паралельно~\cite{nussbaumer1982}.

\textbf{Стадія~2 --- поворотні множники.}
Результати стадії~1 множаться на матрицю поворотних множників:

\begin{equation}
  H_r(k_1) = W_{16}^{r \cdot k_1} \cdot G_r(k_1),
  \quad r, k_1 \in \{0,1,2,3\}.
  \label{eq:stage2}
\end{equation}

При $r = 0$ або $k_1 = 0$ маємо $W_{16}^0 = 1$, тому 7~з 16~множень
є тривіальними. Числові значення 9~нетривіальних множників наведено
у табл.~\ref{tab:twiddle}.

\begin{table}[H]
  \caption{Поворотні множники $W_{16}^{r \cdot k_1}$ (округлено до 3 знаків)}
  \label{tab:twiddle}
  \centering
  \small
  \begin{tabular}{c|cccc}
    \toprule
    $r \backslash k_1$ & 0 & 1 & 2 & 3 \\
    \midrule
    0 & $1$ & $1$ & $1$ & $1$ \\
    1 & $1$ & $0{,}924{-}0{,}383j$ & $0{,}707{-}0{,}707j$ & $0{,}383{-}0{,}924j$ \\
    2 & $1$ & $0{,}707{-}0{,}707j$ & $-j$                  & $-0{,}707{-}0{,}707j$ \\
    3 & $1$ & $0{,}383{-}0{,}924j$ & $-0{,}707{-}0{,}707j$ & $-0{,}924{+}0{,}383j$ \\
    \bottomrule
  \end{tabular}
\end{table}

\textbf{Стадія~3 --- стовпцеві 4-точкові ДПФ.}
Для кожного $k_1$ обчислюється 4-точкове ДПФ стовпця матриці
$\mathbf{H}$:

\begin{equation}
  X(4k_1 + k_2) = \sum_{r=0}^{3} H_r(k_1)\, W_4^{r k_2},
  \quad k_2 \in \{0,1,2,3\}.
  \label{eq:stage3}
\end{equation}

Чотири стовпцевих перетворення~(\ref{eq:stage3}) також є взаємно
незалежними.

\subsubsection*{Оптимізоване 4-точкове ДПФ (метелик radix-4)}

Ядром обох стадій~(\ref{eq:stage1}) і~(\ref{eq:stage3}) є 4-точкове
ДПФ, що завдяки властивостям $W_4 \in \{1, -j, -1, j\}$ реалізується
без нетривіальних множень. Для входів $v_0, v_1, v_2, v_3 \in \mathbb{C}$
вводяться проміжні суми/різниці:

\begin{equation*}
  E_0 = v_0 + v_2, \quad E_1 = v_0 - v_2,
\end{equation*}
\begin{equation*}
  O_0 = v_1 + v_3, \quad O_1 = v_1 - v_3,
\end{equation*}

після чого вихідні відліки отримуються як:

\begin{equation}
  \begin{aligned}
    V_0 &= E_0 + O_0, & V_2 &= E_0 - O_0, \\
    V_1 &= E_1 - j \cdot O_1, & V_3 &= E_1 + j \cdot O_1.
  \end{aligned}
  \label{eq:dft4_butterfly}
\end{equation}

Множення на $j$ у~(\ref{eq:dft4_butterfly}) зводиться до перестановки
дійсної та уявної частин зі зміною знаку
($j \cdot (a + bj) = -b + aj$) і не потребує операції множення у
традиційному сенсі. Таким чином, одне 4-точкове ДПФ потребує
\textbf{8~комплексних додавань і 0 нетривіальних
множень}~\cite{nussbaumer1982}.

\subsubsection*{Підрахунок операцій}

Загальна обчислювальна вартість алгоритму для $N = 16$:

\begin{itemize}
  \item \textbf{Стадія~1:} $4 \cdot 8 = 32$ комплексних додавань, $0$
        нетривіальних множень;
  \item \textbf{Стадія~2:} $9$ нетривіальних комплексних множень
        (що відповідає $36$ дійсним множенням), $16$ комплексних
        додавань;
  \item \textbf{Стадія~3:} $4 \cdot 8 = 32$ комплексних додавань, $0$
        нетривіальних множень.
\end{itemize}

Разом: \textbf{9 нетривіальних комплексних множень} і
\textbf{80 комплексних додавань}, порівняно з $N^2 = 256$ комплексних
множень прямого ДПФ~(\ref{eq:dft}) --- прискорення у
$\approx 28$~разів за кількістю множень~\cite{nussbaumer1982,burrus1985dft}.

\subsubsection*{Формалізація у парадигмі потоку даних}

Алгоритм~(\ref{eq:stage1})--(\ref{eq:stage3}) подається у вигляді
орієнтованого ациклічного графа потоку даних (DFG)
$G = (V, E)$~\cite{dennis1974,veen1986,lee1994}, де:

\begin{itemize}
  \item $V$ --- множина вузлів чотирьох типів
        (\texttt{source}, \texttt{dft4}, \texttt{twiddle}, \texttt{sink});
  \item $E \subseteq V \times V$ --- множина спрямованих дуг, що
        є FIFO-чергами токенів.
\end{itemize}

\textbf{Правило активації.} Вузол $v \in V$ є готовим до виконання
тоді і лише тоді, коли на кожному з його вхідних каналів присутній
щонайменше один токен. При активації вузол поглинає вхідні токени,
виконує обчислення, визначене його типом, та породжує відповідні
вихідні токени, що розміщуються на вихідних дугах.

\textbf{Часова модель.} Вузол має затримку виконання \texttt{latency},
а дуга --- затримку доставки \texttt{delay}. Токен, породжений у
момент $t$, стає готовим до споживання у момент $t + \texttt{delay}$.
Логічний час $\tau$ симулятора інкрементується дискретно на кроці
$\mathrm{d}\tau$, що регулюється множником швидкості.

\subsubsection*{Псевдокод планувальника}

Планувальник (ядро \texttt{DataflowRuntime}) на кожному кроці
симуляції виконує дві дії: доставку токенів із FIFO-черг і запуск
готових вузлів. Узагальнений псевдокод:

\begin{figure}[H]
\begin{verbatim}
procedure TICK(dt, maxFires)
    tau <- tau + dt
    for each edge e in E
        while e.queue.head.t + e.delay <= tau
            token <- e.queue.pop()
            v <- target(e); port <- targetPort(e)
            v.inputs[port] <- token
    fires <- 0
    for each node v in V  order by id
        if v is ready and not busy and fires < maxFires
            emit onFire(v, tau)
            v.busy <- tau + v.latency
            outs <- compute(v.kind, v.inputs)
            schedule outs on outgoing edges at time v.busy
            fires <- fires + 1
    for each node v in V
        if v.busy <= tau then v.busy <- 0
\end{verbatim}
\caption{Псевдокод однієї ітерації ядра симуляції \texttt{DataflowRuntime.tick()}.
         Джерело: розроблено автором.}
\label{fig:tick_pseudo}
\end{figure}

Ця процедура забезпечує детермінованість (NF5): при фіксованому порядку
обходу вузлів результат симуляції не залежить від реального часу
виконання у браузері.


% ─────────────────────────────────────────────────
\subsection{Загальна архітектура системи}

Архітектура симулятора подається на двох рівнях моделі C4~\cite{c4model}
--- Context і Container, --- доповнених діаграмою потоків даних (DFD)
для внутрішньої логіки симуляції.

\subsubsection*{Контекстний рівень (C4 Level 1)}

На контекстному рівні система представляється як «чорна скринька»,
що взаємодіє з одним актором --- \emph{користувачем} (студент або
викладач), та з двома зовнішніми сутностями: \emph{браузером} (runtime-
середовище) і \emph{локальним сховищем браузера} (для збереження
пресетів вхідних сигналів і параметрів). Застосунок завантажується
одноразово зі статичного веб-хостингу та надалі працює автономно.

\subsubsection*{Рівень контейнерів (C4 Level 2)}

На рівні контейнерів симулятор декомпозується на чотири логічні
підсистеми, що виконуються у межах одного процесу браузера:

\begin{enumerate}
  \item \textbf{Simulation Core} --- ядро симуляції (\texttt{DataflowRuntime},
        \texttt{scheduler}, \texttt{math}): реалізує псевдокод
        рис.~\ref{fig:tick_pseudo}, зберігає стан графа та черг;
  \item \textbf{Graph Generator} --- модуль побудови DFG для заданого
        $N$ та параметрів затримок;
  \item \textbf{State Store} --- централізоване сховище стану, до
        якого підписуються всі UI-компоненти;
  \item \textbf{UI Layer} --- набір React-компонентів:
        \texttt{GraphView} (полотно React Flow), \texttt{SpectrumView}
        (діаграма $|X(k)|$), \texttt{Metrics} (пропускна здатність,
        затримка, черги), \texttt{Controls} (пуск/пауза/швидкість).
\end{enumerate}

Взаємодії між контейнерами: UI Layer читає стан зі State Store та
викликає команди (пуск, пауза); State Store делегує обробку команд
до Simulation Core; Simulation Core публікує події
(\texttt{onToken}, \texttt{onFire}, \texttt{onOutput}), на які
State Store оновлює відповідні поля; Graph Generator викликається
один раз при ініціалізації та за кожної зміни $N$ чи latency.

\subsubsection*{Граф потоку даних обчислювача}

Детальна структура DFG для $N = 16$ складається з \textbf{56~вузлів}
і \textbf{64~дуг}, розподілених за п'ятьма ярусами:

\begin{enumerate}
  \item 16~вузлів \texttt{source} (src$_0 \ldots$src$_{15}$) ---
        генерують токени вхідних відліків $x(n)$;
  \item 4~вузли \texttt{dft4} (row-0 \ldots row-3) --- виконують
        стадію~1~(\ref{eq:stage1});
  \item 16~вузлів \texttt{twiddle} (tw-$r$-$k_1$) --- виконують
        стадію~2~(\ref{eq:stage2});
  \item 4~вузли \texttt{dft4} (col-0 \ldots col-3) --- виконують
        стадію~3~(\ref{eq:stage3});
  \item 16~вузлів \texttt{sink} (snk$_0 \ldots$snk$_{15}$) ---
        збирають спектральні відліки $X(k)$.
\end{enumerate}

Характеристики вузлів і дуг зведено у
табл.~\ref{tab:nodes}--\ref{tab:edges}.

\begin{table}[H]
  \caption{Типи вузлів DFG та їхні характеристики}
  \label{tab:nodes}
  \centering
  \small
  \begin{tabular}{lccccl}
    \toprule
    Тип вузла & К-сть & Вх. портів & Вих. портів & Затримка & Функція \\
    \midrule
    \texttt{source}  & 16 & 0 & 1 & ---  & Генерація $x(n)$ \\
    \texttt{dft4}    &  8 & 4 & 4 & 400  & 4-точкове ДПФ~(\ref{eq:dft4_butterfly}) \\
    \texttt{twiddle} & 16 & 1 & 1 & 200  & Множення на $W_{16}^{r k_1}$ \\
    \texttt{sink}    & 16 & 1 & 0 & ---  & Збір $X(k)$ \\
    \midrule
    \multicolumn{2}{l}{\textbf{Разом: 56 вузлів}} & & & & \\
    \bottomrule
  \end{tabular}
\end{table}

\begin{table}[H]
  \caption{Типи дуг DFG}
  \label{tab:edges}
  \centering
  \small
  \begin{tabular}{llcc}
    \toprule
    З вузла & До вузла & К-сть дуг & Затримка \\
    \midrule
    \texttt{source}      & \texttt{dft4} (рядковий)  & 16 & 600 \\
    \texttt{dft4} (ряд.) & \texttt{twiddle}          & 16 & 600 \\
    \texttt{twiddle}     & \texttt{dft4} (стовпц.)   & 16 & 600 \\
    \texttt{dft4} (ст.)  & \texttt{sink}             & 16 & 600 \\
    \midrule
    \multicolumn{2}{l}{\textbf{Разом: 64 дуги}} & & \\
    \bottomrule
  \end{tabular}
\end{table}

\subsubsection*{Критичний шлях і коефіцієнт паралелізму}

Критичний шлях у DFG (найдовший шлях від \texttt{source} до
\texttt{sink} з урахуванням затримок вузлів і дуг) становить:

\begin{equation}
  T_{\mathrm{kr}} = 400 + 600 + 200 + 600 + 400 = 2200
  \label{eq:critical_path}
\end{equation}

умовних одиниць часу. Послідовний час виконання (при одному
активному вузлі) становив би:

\begin{equation}
  T_{\mathrm{ps}} = 8 \cdot 400 + 16 \cdot 200 = 6400.
  \label{eq:seq_time}
\end{equation}

Коефіцієнт паралелізму:

\begin{equation}
  K_p = \frac{T_{\mathrm{ps}}}{T_{\mathrm{kr}}} = \frac{6400}{2200} \approx 2{,}9,
  \label{eq:parallelism}
\end{equation}

тобто паралельне виконання незалежних вузлів забезпечує майже
трикратне прискорення порівняно з послідовним обчисленням~\cite{veen1986}.

\begin{figure}[H]
  \centering
  \IfFileExists{images/dfg.png}{%
    \includegraphics[width=0.95\textwidth]{images/dfg.png}%
  }{%
    \fbox{\parbox[c][0.28\textheight][c]{0.95\textwidth}{\centering
      Рисунок буде додано після генерації у симуляторі.\\
      Файл: \texttt{tex/images/dfg.png}}}%
  }
  \caption{Граф потоку даних обчислювача 16-точкового ШПФ radix-4
           $4 \times 4$. Джерело: розроблено автором.}
  \label{fig:dfg}
\end{figure}


% ─────────────────────────────────────────────────
\subsection{Структури даних}

Інформаційне забезпечення симулятора визначається набором типів
даних, що є програмними відповідниками формальних об'єктів DFG
з підрозд.~2.2.

\subsubsection*{Комплексні числа}

Усі числові значення, з якими оперує симулятор, є комплексними.
Вони подаються структурою \texttt{Complex} у форматі
\{\,\texttt{re}: number,\ \texttt{im}: number\,\}, де \texttt{re},
\texttt{im} --- дійсна та уявна частини числа
$z = \texttt{re} + j \cdot \texttt{im}$, збережені у подвійній
точності IEEE~754 (64~біти, машинна точність
$\varepsilon_{\mathrm{mach}} \approx 2.22 \cdot 10^{-16}$).

\subsubsection*{Тип \texttt{Token}}

Тип \texttt{Token} є програмним відповідником токена DFG, що
переміщується по дугам графа (табл.~\ref{tab:token}).

\begin{table}[H]
  \caption{Поля типу \texttt{Token}}
  \label{tab:token}
  \centering
  \small
  \begin{tabular}{lll}
    \toprule
    Поле & Тип & Призначення \\
    \midrule
    \texttt{id}      & \texttt{string}               & Унікальний ідентифікатор токена \\
    \texttt{value}   & \texttt{Complex}              & Комплексне значення токена \\
    \texttt{t}       & \texttt{number}               & Логічний час надходження на дугу \\
    \texttt{originT} & \texttt{number}               & Час народження у \texttt{source} \\
    \bottomrule
  \end{tabular}
\end{table}

Поле \texttt{t} визначає момент готовності токена: він стає доступним
вузлу-приймачу в момент $t + \texttt{edge.delay}$. Поле
\texttt{originT} зберігається незмінним упродовж усього маршруту
токена і використовується для вимірювання наскрізної
затримки~(\ref{eq:critical_path}).

\subsubsection*{Тип \texttt{Edge}}

Тип \texttt{Edge} описує спрямований канал між двома портами
вузлів (табл.~\ref{tab:edge}).

\begin{table}[H]
  \caption{Поля типу \texttt{Edge}}
  \label{tab:edge}
  \centering
  \small
  \begin{tabular}{lll}
    \toprule
    Поле & Тип & Призначення \\
    \midrule
    \texttt{id}     & \texttt{string}                  & Унікальний ідентифікатор дуги \\
    \texttt{from}   & \texttt{\{node, port\}}          & Вузол і порт-джерело \\
    \texttt{to}     & \texttt{\{node, port\}}          & Вузол і порт-приймач \\
    \texttt{delay}  & \texttt{number}                  & Затримка доставки токена (600) \\
    \texttt{label}  & \texttt{string?}                 & Мітка на ребрі ($W_{16}^{rk_1}$) \\
    \texttt{labelValue} & \texttt{Complex?}            & Числове значення поворотного множника \\
    \bottomrule
  \end{tabular}
\end{table}

\subsubsection*{Тип \texttt{NodeSpec} і \texttt{Graph}}

Тип \texttt{NodeSpec} визначає специфікацію вузла DFG
(табл.~\ref{tab:nodespec}), а тип \texttt{Graph = \{nodes:
NodeSpec[], edges: Edge[]\}} --- повний граф обчислювача. Для
$N = 16$ об'єкт \texttt{Graph} містить 56~специфікацій вузлів та
64~дуги.

\begin{table}[H]
  \caption{Поля типу \texttt{NodeSpec}}
  \label{tab:nodespec}
  \centering
  \small
  \begin{tabular}{lll}
    \toprule
    Поле & Тип & Призначення \\
    \midrule
    \texttt{id}       & \texttt{string}    & Ідентифікатор вузла \\
    \texttt{kind}     & \texttt{NodeKind}  & Тип (source, dft4, twiddle, sink) \\
    \texttt{inPorts}  & \texttt{string[]}  & Список вхідних портів \\
    \texttt{outPorts} & \texttt{string[]}  & Список вихідних портів \\
    \texttt{latency}  & \texttt{number}    & Тривалість обчислення (400 або 200) \\
    \texttt{params}   & \texttt{object?}   & Параметри (для twiddle: $N, k$) \\
    \bottomrule
  \end{tabular}
\end{table}

Тип \texttt{NodeKind} є переліковим: \texttt{'source' | 'dft4' |
'twiddle' | 'sink'}. Він використовується як ключ у шаблоні
Strategy (підрозд.~2.1) для вибору функції обчислення вузла.

\subsubsection*{Проміжні матриці обчислення}

Алгоритм~(\ref{eq:stage1})--(\ref{eq:stage3}) оперує чотирма
матрицями даних розміром $4 \times 4$:

\begin{align}
  \mathbf{A}[r][s]   &= x(4s + r),            && \text{вхідні дані,}  \label{eq:A_def} \\
  \mathbf{G}[r][k_1] &= \mathrm{DFT4}(\mathbf{A}[r][\cdot])[k_1],
                     && \text{після стадії 1,}  \label{eq:G_def} \\
  \mathbf{H}[r][k_1] &= W_{16}^{r \cdot k_1} \cdot \mathbf{G}[r][k_1],
                     && \text{після стадії 2,}  \label{eq:H_def} \\
  X(4k_1 + k_2)      &= \mathrm{DFT4}(\mathbf{H}[\cdot][k_1])[k_2].
                     && \text{вихідний спектр.} \label{eq:X_def}
\end{align}

\subsubsection*{Числовий приклад: $x(n) = n$}

Для контрольної верифікації використано лінійний вхідний сигнал
$x(n) = n$, $n = 0 \ldots 15$. Проміжні значення для першого рядка
($r = 0$) і перших чотирьох спектральних відліків наведено у
табл.~\ref{tab:numerical}.

\begin{table}[H]
  \caption{Проміжні значення алгоритму для $x(n) = n$, рядок $r = 0$}
  \label{tab:numerical}
  \centering
  \small
  \begin{tabular}{lcccc}
    \toprule
    Матриця & Індекс 0 & Індекс 1 & Індекс 2 & Індекс 3 \\
    \midrule
    $\mathbf{A}[0][\cdot]$   & $0$       & $4$              & $8$         & $12$ \\
    $\mathbf{G}[0][\cdot]$   & $24{+}0j$ & $-8{+}8j$        & $-8{+}0j$   & $-8{-}8j$ \\
    $\mathbf{H}[0][\cdot]$   & $24{+}0j$ & $-8{+}8j$        & $-8{+}0j$   & $-8{-}8j$ \\
    $X(k)$, $k{=}0{\ldots}3$ & $120{+}0j$ & ${\approx}-8{+}40j$ & ${\approx}-8{+}19j$ & ${\approx}-8{+}12j$ \\
    \bottomrule
  \end{tabular}
\end{table}

Рядок $\mathbf{H}[0][\cdot]$ збігається з $\mathbf{G}[0][\cdot]$,
оскільки $W_{16}^{0 \cdot k_1} = 1$ для всіх $k_1$. Точне значення
$X(0) = \sum_{n=0}^{15} n = 120$ підтверджує коректність стадії~3.


% ─────────────────────────────────────────────────
\subsection{Механізми обміну даними та інтеграції}

Оскільки симулятор не має серверної частини, класичні механізми
мережевої інтеграції (REST API, WebSocket, gRPC) у ньому не
застосовуються. Замість них координація підсистем реалізується
через \emph{внутрішній подійний канал} і \emph{реактивне сховище
стану}, описані нижче.

\subsubsection*{Шина подій симуляції}

Ядро \texttt{DataflowRuntime} публікує три типи подій, що утворюють
внутрішній API симулятора (табл.~\ref{tab:events}).

\begin{table}[H]
  \caption{Події шини симуляції}
  \label{tab:events}
  \centering
  \small
  \begin{tabular}{p{2.2cm}p{5.2cm}p{5.8cm}}
    \toprule
    Подія & Умова виникнення & Підписники \\
    \midrule
    \texttt{onToken}  & Токен надійшов на дугу $e \in E$ & Heatmap черг, Metrics (queue size) \\
    \texttt{onFire}   & Активовано вузол $v \in V$ & Heatmap активності вузлів~(\ref{eq:heat_decay}) \\
    \texttt{onOutput} & Токен надійшов до \texttt{sink} & SpectrumView, Metrics (throughput, latency) \\
    \bottomrule
  \end{tabular}
\end{table}

Події реалізують шаблон Publisher/Subscriber (підрозд.~2.1):
видавець (ядро) не має прямих посилань на підписників, а лише
передає об'єкт події у список зареєстрованих обробників. Це
забезпечує ізоляцію бізнес-логіки від UI та дозволяє додавати нові
підписники (наприклад, експорт event~log у файл) без модифікації
ядра.

\subsubsection*{Реактивне сховище стану}

Стан симулятора централізований у сховищі \texttt{simStore}, що
містить такі основні групи полів:

\begin{itemize}
  \item \textit{Стан графа}: \texttt{graph: Graph},
        \texttt{queues: Map<EdgeId, Token[]>},
        \texttt{nodeBusy: Map<NodeId, number>};
  \item \textit{Стан виконання}: \texttt{status: 'IDLE' | 'RUNNING' |
        'PAUSED' | 'DONE'}, \texttt{simTime: number},
        \texttt{speed: number};
  \item \textit{Метрики}: \texttt{throughputEMA},
        \texttt{latencyEMA}, \texttt{queueSizeEMA},
        \texttt{activity: Map<NodeId, number>};
  \item \textit{Результат}: \texttt{spectrum: Complex[16]},
        \texttt{eventLog: Event[]}.
\end{itemize}

UI-компоненти підписуються на окремі поля сховища за допомогою
селекторів (pattern Selector). Компонент перемальовується лише
при зміні саме тих полів, на які він підписаний, що забезпечує
стабільні 60~fps (NF3) навіть при інтенсивному потоці подій від
ядра.

\subsubsection*{Механізм пресетів та збереження стану}

Для підтримки повторюваних демонстраційних сценаріїв (UC1 з
підрозд.~1.5) передбачено механізм пресетів. Пресет --- це JSON-об'єкт,
що містить: вхідний сигнал $\{x_n\}$, значення \texttt{latency} і
\texttt{delay}, початкову швидкість симуляції. Пресети зберігаються
у Local Storage браузера за ключем \texttt{simulator.presets} і
завантажуються через окремий компонент \texttt{PresetSelector}.
Експорт пресету у файл (формат \texttt{.json}) реалізується
стандартним Web API \texttt{Blob + URL.createObjectURL}.

Експорт event~log (UC9) реалізований аналогічно: журнал подій
серіалізується у JSON або CSV і завантажується користувачем без
передачі на сервер, що відповідає принципу конфіденційності
(підрозд.~1.4).


% ─────────────────────────────────────────────────
\subsection{Інтерфейс користувача та взаємодія із системою}

Інтерфейс симулятора спроєктовано за принципом \emph{інформаційної
ієрархії}: найбільша площа екрана відводиться під центральний
об'єкт дослідження --- граф потоку даних, --- а допоміжні панелі
(метрики, спектр, керування) розміщуються по периметру, не
перекриваючи граф.

\subsubsection*{Структура інтерфейсу}

Головний екран застосунку розділений на чотири функціональні зони:

\begin{enumerate}
  \item \textbf{Центральна зона (GraphView)} --- полотно React~Flow
        з вузлами та анімованими дугами. Користувач може перетягувати
        вузли, змінювати масштаб колесом миші, клікати по вузлу для
        перегляду деталей (поточні значення на входах/виходах).
  \item \textbf{Верхня панель (Controls)} --- кнопки «Пуск», «Пауза»,
        «Крок», «Скидання», повзунок множника швидкості ($0.1\times$
        до $10\times$), селектор пресету вхідного сигналу, поле
        редагування \texttt{latency}/\texttt{delay}.
  \item \textbf{Права панель (SpectrumView + Metrics)} --- стовпчикова
        діаграма $|X(k)|$ у лінійному або логарифмічному масштабі та
        три числових індикатори (throughput, latency, queue size).
  \item \textbf{Нижня панель (TerminalOutput)} --- журнал подій у
        реальному часі з можливістю фільтрації за типом події
        (\texttt{onToken}/\texttt{onFire}/\texttt{onOutput}) і
        експорту у файл.
\end{enumerate}

\subsubsection*{UX-логіка взаємодії}

Базовий user-flow для студента (UC3 з підрозд.~1.5) складається з
трьох кроків, доступних у межах двох кліків від головного екрана
(NF8):

\begin{enumerate}
  \item Вибір пресету вхідного сигналу (синусоїда, імпульс, шум)
        або введення власних відліків через компонент
        \texttt{InputEditor};
  \item Натискання кнопки «Пуск» --- ядро починає викликати
        \texttt{tick()} у циклі \texttt{requestAnimationFrame}
        з кроком $\mathrm{d}\tau = 16.67$~мс $\cdot$ speed;
  \item Спостереження за рухом токенів у GraphView, зміною метрик у
        Metrics та поступовим заповненням SpectrumView.
\end{enumerate}

У режимі паузи доступний покроковий режим \emph{Step}, у якому одна
кнопка викликає один виклик \texttt{tick()}: це дозволяє детально
дослідити виконання алгоритму у навчальних цілях.

\subsubsection*{Візуальні акценти на потоковій природі алгоритму}

Для підсилення сприйняття потокової природи алгоритму UI використовує
три допоміжні візуальні механізми:

\begin{itemize}
  \item \textbf{Анімація токенів.} Кожен токен відображається як
        кружечок, що переміщується по дузі в інтервалі часу
        $[t, t+\texttt{delay}]$ з плавною інтерполяцією позиції.
        На кружечку виводиться числове значення (дійсна та уявна
        частини) для прозорого спостереження обчислень.
  \item \textbf{Теплова карта активності вузлів.} Заливка кожного
        вузла модулюється значенням $a_i(\tau)$, що оновлюється за
        експоненціальним законом:
        \begin{equation}
          a_i(\tau + \mathrm{d}\tau) = \alpha \cdot a_i(\tau) + \Delta a_i,
          \quad \alpha = 0{,}92,
          \label{eq:heat_decay}
        \end{equation}
        де $\Delta a_i = 1$ при подіях \texttt{onFire} вузла $i$ та
        $\Delta a_i = 0$ в інших випадках. Вузли, що нещодавно
        часто активувалися, відображаються яскравішими («гарячими»).
  \item \textbf{Стан busy на вузлах.} Поки вузол перебуває у стані
        виконання ($v.\text{busy} > \tau$), навколо нього
        відображається кільцевий індикатор прогресу, що ілюструє
        затримку \texttt{latency} вузла.
\end{itemize}

\subsubsection*{Макет інтерфейсу}

Загальний вигляд інтерфейсу подано на рис.~\ref{fig:ui}.

\begin{figure}[H]
  \centering
  \IfFileExists{images/simulator_ui.png}{%
    \includegraphics[width=0.95\textwidth]{images/simulator_ui.png}%
  }{%
    \fbox{\parbox[c][0.24\textheight][c]{0.95\textwidth}{\centering
      Скріншот буде додано після фінального тестування.\\
      Файл: \texttt{tex/images/simulator\_ui.png}}}%
  }
  \caption{Макет інтерфейсу веб-симулятора потокового 16-точкового
           ШПФ: центральна зона GraphView, верхня панель Controls,
           права панель SpectrumView + Metrics, нижня панель
           TerminalOutput. Джерело: розроблено автором.}
  \label{fig:ui}
\end{figure}


% ─────────────────────────────────────────────────
\subsection*{Висновки до розділу~2}

У другому розділі спроєктовано програмне рішення, що відповідає
вимогам, сформульованим у підрозд.~1.5.

\emph{Обрано архітектурний підхід} клієнтського SPA без серверної
частини, що забезпечує одночасне виконання вимог NF2 (відгук
$\le 100$~мс), NF4 (портативність Chrome/Firefox/Edge), NF7 (робота
без встановлення) і NF5 (детермінованість). Вибір реактивного
рендерингу та централізованого сховища стану обґрунтований як
компроміс між продуктивністю та супроводжуваністю коду. Застосовано
шаблони проєктування Observer, State, Strategy, Factory і Facade;
зв'язок між вимогами та рішеннями показано у
табл.~\ref{tab:req_to_tech}.

\emph{Побудовано математичну модель} алгоритму ШПФ radix-4 для
$N = 16$ у вигляді тристадійної $4 \times 4$ декомпозиції
(\ref{eq:stage1})--(\ref{eq:stage3}), що потребує 9~нетривіальних
комплексних множень і 80~комплексних додавань проти 256~множень
прямого ДПФ. Сформульовано формальну модель потоку даних із
правилом активації та часовою семантикою, подано псевдокод
планувальника (рис.~\ref{fig:tick_pseudo}).

\emph{Спроєктовано архітектуру системи} на рівнях C4 Context і
Container з декомпозицією на чотири підсистеми (Simulation Core,
Graph Generator, State Store, UI Layer). Деталізовано граф
обчислювача: 56~вузлів чотирьох типів і 64~дуги з однаковою
затримкою 600~одиниць, критичним шляхом $T_{\mathrm{kr}} = 2200$
і коефіцієнтом паралелізму $K_p \approx 2{,}9$.

\emph{Визначено структури даних:} \texttt{Token}, \texttt{Edge},
\texttt{NodeSpec}, \texttt{Graph}, що є програмними відповідниками
формальних об'єктів DFG. Комплексні числа подаються у форматі
IEEE~754 double, що гарантує запас точності відносно
$\varepsilon = 10^{-9}$ (NF1).

\emph{Обрано механізми обміну даними:} внутрішня шина подій
(\texttt{onToken}, \texttt{onFire}, \texttt{onOutput}) для ізоляції
UI від ядра, реактивне сховище стану та пресети у Local~Storage
для повторюваних сценаріїв демонстрації.

\emph{Спроєктовано інтерфейс користувача} з інформаційною ієрархією
(центральний GraphView, периферійні Controls, SpectrumView, Metrics,
TerminalOutput), трьома візуальними механізмами для підсилення
потокової природи алгоритму (анімація токенів, теплова карта
активності, індикатор busy) та базовим user-flow у межах двох
кліків (NF8).

Описані проєктні рішення утворюють повне технічне завдання для
реалізації, що подається у розділі~3.